\documentclass[11pt,a4paper]{article}

\usepackage[T1]{fontenc}
\usepackage[margin=2.5cm]{geometry}
\usepackage{booktabs}
\usepackage{array}
\usepackage{enumitem}
\usepackage[hidelinks]{hyperref}

\setlength{\parskip}{0.5em}
\setlength{\parindent}{0pt}
\setlist{itemsep=0.1em,topsep=0.3em}

\newcommand{\code}[1]{\texttt{#1}}
% Marks text the author still has to fill in; remove before submission.
\newcommand{\todo}[1]{\textbf{[TODO: #1]}}

\title{Predictive Analysis of AI-Bot Evolution Using NLP-Based Canon Event Extraction\\[0.3em]
\large Introduction and System Description}
\author{Project team}
\date{}

\begin{document}
\maketitle

%==============================================================================
\section{Introduction}
\label{sec:intro}

\subsection{Motivation and research question}
AI systems have moved quickly from language models that only generate text to systems that call tools, plan
multi-step tasks, operate a computer and coordinate with other agents. Evidence of this change is spread over a
large and fast-moving collection of research papers, laboratory blog posts, release notes and incident reports.
These sources are unstructured, repetitive and scattered over time, so it is hard to say when a capability
first appeared, how it spread, or whether the overall pattern of development is changing.

This project asks: \emph{can NLP-based extraction of significant AI-agent developments reveal a measurable
evolution timeline, identify transitions between phases of AI-agent development, and support a data-driven
forecast of an emerging next phase?}

We deliberately do not claim to predict a date for an ``AI revolution''. The measurable goal is narrower: to
forecast an explicitly defined quantity, the dominant event category in the next month, from patterns in the
historical event sequence.

\subsection{Approach}
The central object of the project is the \emph{AI canon event}: a significant, dated development in a
language-model-based or agentic AI system that represents a meaningful change in capability, autonomy,
interaction, deployment or ecosystem behaviour. Each event is stored as a structured record (organisation,
system, capability, action, autonomy level, interaction type, impact, source and confidence) and assigned to one
of eight categories: capability breakthrough, agentic development, multi-agent coordination, autonomous action,
unexpected behaviour or incident, infrastructure and ecosystem, adoption and scale, and research milestone.
The categories and their rules are fixed in a single configuration file shared by every component.

The processing chain is: (1)~collect dated documents; (2)~clean and deduplicate them; (3)~extract events with
NER, temporal extraction and event-trigger detection; (4)~classify and normalise events and link reports of the
same real event; (5)~place events on a timeline; (6)~compute monthly features; (7)~detect trends and phase
changes; (8)~forecast the dominant category of the next month using only past data; and (9)~present the result
with its supporting events and sources. The conceptual phase sequence from the project brief (language models,
tool use, agents, autonomous agents, multi-agent systems, emergent behaviour, large-scale ecosystems) is used as
an initial hypothesis that the data may confirm or contradict, not as an assumed fact.

\subsection{Scope of this part of the report}
The project was carried out by four members who all took part in annotation, review and evaluation. This part
describes the \emph{system}: how the corpus was collected and cleaned (Section~\ref{sec:collection} and
\ref{sec:cleaning}), how documents and events are stored and queried (Section~\ref{sec:db}), the interactive
application (Section~\ref{sec:app}) and how the work was kept reproducible (Section~\ref{sec:repro}). Section~%
\ref{sec:limits} lists the limits of the data infrastructure. Event extraction, linking, phase analysis,
forecasting and their evaluation are described in the other parts of the report, and the system described here
is built to display their outputs.

\subsection{A note on interpretation}
News and reports are an imperfect proxy for what AI systems can actually do. Media attention, source choice,
duplicate reporting, missing events and changing terminology all distort a timeline built from text. All
results in this report are therefore analyses of \emph{reported} developments in our corpus. The application
keeps every number traceable to its source documents so that readers can check it.

%==============================================================================
\section{System Description}
\label{sec:system}

\subsection{Overview}
The system is a Python code base (version~3.13) organised as a package, \code{ainevo}, plus command-line
scripts and a Streamlit application. Table~\ref{tab:overview} summarises the components.

\begin{table}[ht]
\centering
\small
\begin{tabular}{@{}p{3.3cm}p{11.7cm}@{}}
\toprule
Component & Purpose \\
\midrule
Collection & Downloads dated documents from arXiv, laboratory blogs and GitHub release pages \\
Preprocessing & Merges sources, cleans text and dates, removes duplicates, applies monthly caps \\
Storage & SQLite database holding documents, events and event--document links; full-text search \\
Later stages & Extraction, linking, features, analysis and forecasting; described in the other parts of the report \\
Application & Streamlit pages for search, timeline, phases, forecast and evidence \\
Configuration & YAML files for sources, the event taxonomy and the forecasting split \\
\bottomrule
\end{tabular}
\caption{Components of the system.}
\label{tab:overview}
\end{table}

Raw data, the cleaned corpus and the database are generated files and are not kept in version control. Every
step can be re-run from the command line (Section~\ref{sec:repro}).

%------------------------------------------------------------------------------
\subsection{Data collection}
\label{sec:collection}
The corpus combines three document families, chosen because they complement each other: papers show research
milestones, laboratory posts give precisely dated releases and capability announcements, and release notes of
agent frameworks show how tooling and infrastructure change.

\begin{description}[leftmargin=0.5cm,style=nextline]
\item[Research papers.] Titles, abstracts, dates and links are retrieved from the arXiv API with category and
keyword filters (for example agent, tool use, planning, memory, multi-agent, computer use). A per-year target
keeps early years from being empty while the volume of agent papers grows sharply after 2022.
\item[Laboratory and company blogs.] Posts from OpenAI, Anthropic, Google DeepMind, Google Research and
Hugging Face are collected from their feeds, and the full page text is fetched where possible. The \code{text\_source}
field records whether the text is a full page or only a feed summary. Before fetching pages the collector checks
\code{robots.txt}.
\item[Release notes.] GitHub releases of four agent frameworks (LangChain, CrewAI, AutoGPT and AutoGen) are
collected through the GitHub API; an optional access token raises the rate limit.
\end{description}

All collectors follow the same rules: respect \code{robots.txt}, terms of service and rate limits; record the
retrieval date and the query or rule that found each document; keep the original source information; and never
guess a missing publication date. Each document receives a stable identifier,
\code{doc\_id} = first two letters of the source type + ``\_'' + the first ten characters of the MD5 hash of its URL, so
the same page always has the same identifier.

\paragraph{Document record.} Every document has the fields \code{doc\_\allowbreak{}id}, \code{title}, \code{publication\_\allowbreak{}date},
\code{source}, \code{source\_\allowbreak{}type}, \code{url}, \code{text}, \code{retrieval\_\allowbreak{}date}, \code{query\_\allowbreak{}used} and a
\code{metadata} object (flags \code{llm\_\allowbreak{}terms\_\allowbreak{}found}, \code{low\_\allowbreak{}signal\_\allowbreak{}release} and \code{text\_\allowbreak{}source}).

%------------------------------------------------------------------------------
\subsection{Cleaning, deduplication and the frozen corpus}
\label{sec:cleaning}
The raw files are merged and cleaned in one pipeline. It removes boilerplate and irrelevant text, normalises
dates, removes exact duplicates and removes near-duplicates using MinHash similarity, while preserving the
original source information. \todo{insert document counts before and after cleaning from cleaning\_report.json}

Two design choices affect later analysis. First, \emph{monthly caps} limit how many documents one source can
contribute (OpenAI 15 per month, Hugging Face 8 per month, release notes 8 per repository per month). Without
them a few prolific sources would dominate the monthly event counts that the forecast uses. Second, documents
that were hand-annotated are never removed by cleaning, so the gold labels always refer to documents that exist
in the final corpus.

The corpus was frozen with a fixed data cutoff so that the forecasting test set cannot drift. The frozen file
\code{corpus\_v1\_clean.jsonl} contains 3{,}018 documents from 2017 to September 2026 and is identified by its
SHA-256 value
\begin{center}\small\code{48d3d61e15963f13a78a3b7e030d7882fa9b4a75e3ce3d4b78d2560208a88b09.}\end{center}
The last complete month is
September 2026; the 29 documents dated October 2026 stay in the corpus but are outside the forecasting series.
Table~\ref{tab:corpus} gives the composition.

\begin{table}[ht]
\centering
\small
\begin{tabular}{@{}lrrrr@{}}
\toprule
Year & Lab blogs & Papers & Release notes & Total \\
\midrule
2017 & 27 & 41 & 0 & 68 \\
2018 & 18 & 42 & 0 & 60 \\
2019 & 18 & 41 & 0 & 59 \\
2020 & 16 & 42 & 0 & 58 \\
2021 & 43 & 42 & 0 & 85 \\
2022 & 92 & 82 & 0 & 174 \\
2023 & 161 & 153 & 135 & 449 \\
2024 & 225 & 192 & 176 & 593 \\
2025 & 337 & 195 & 222 & 754 \\
2026 & 371 & 154 & 193 & 718 \\
\midrule
Total & 1{,}308 & 984 & 726 & 3{,}018 \\
\bottomrule
\end{tabular}
\caption{Frozen corpus by year and source type. Release notes start in 2023 because the collected framework
releases begin then. The 2026 row includes the 29 documents dated October 2026.}
\label{tab:corpus}
\end{table}

By source the corpus holds arXiv 984, OpenAI 613, Hugging Face 434, LangChain 341, CrewAI 175, Anthropic 149,
AutoGPT 115, AutoGen 95, Google DeepMind 61 and Google Research 51 documents.

%------------------------------------------------------------------------------
\subsection{Event database}
\label{sec:db}
Documents and events are stored in a single SQLite file built from scratch by one script, so the same inputs
always give the same database. SQLite was chosen because the project is small, runs on every member's laptop
and needs no server.

\paragraph{Tables.}
\begin{itemize}
\item \code{documents}: the document record above, plus a year column and indexes on date, type and year. A
full-text index (SQLite FTS5 with Porter stemming) covers titles and text.
\item \code{events}: one row per canon event, with the event fields of the project schema, the month used by
the monthly features, the \code{origin} of the event and the classifier split used for evaluation.
\item \code{event\_sources}: links each event to every document that reports it, so an event can have several
sources after linking.
\end{itemize}

\paragraph{Origins and validation.} Every event has an origin: \code{gold} (labelled by hand by the team under
the frozen guidelines), \code{extracted} (output of the automatic extraction pipeline) or \code{seed}. Before an
event is inserted it is checked against the taxonomy file: category, capability tag, interaction type, autonomy
level (0--4 or not applicable), phase hint (1--7), date format, and the confidence scale, which is 1--3 for
annotator confidence on gold events and a probability between 0 and 1 for extracted events. An invalid row is
rejected and reported without blocking the others.

\paragraph{Gold merges.} The team ruled that two gold annotations describe the same real intrusion at one
organisation in July 2026. The loader applies such rulings from a small reference table, located by annotation
identifier because event identifiers are renumbered between versions, keeps one event and adds the second
document to its sources. The database therefore holds 216 gold events (from 217 annotated events), the same
set that the forecasting features are built from.

\paragraph{Contents.} Table~\ref{tab:events} lists the final database.

\begin{table}[ht]
\centering
\small
\begin{tabular}{@{}lrr@{}}
\toprule
Category & Gold & Extracted \\
\midrule
Multi-agent coordination & 57 & 249 \\
Research milestone & 39 & 140 \\
Agentic development & 30 & 19 \\
Infrastructure and ecosystem & 28 & 131 \\
Capability breakthrough & 28 & 222 \\
Adoption and scale & 15 & 85 \\
Autonomous action & 11 & 2 \\
Unexpected behaviour or incident & 8 & 23 \\
\midrule
Total & 216 & 871 \\
\bottomrule
\end{tabular}
\caption{Events in the database by category and origin. The extracted events are provisional model output
(Section~\ref{sec:limits}).}
\label{tab:events}
\end{table}

\paragraph{Query layer.} The storage module offers document search (free text plus filters on source, source
type and year) and event queries filtered by category, organisation, capability, autonomy level, month range and
origin. Each event returned carries its list of source documents, which is what makes every displayed result
traceable.

%------------------------------------------------------------------------------
\subsection{Interactive application}
\label{sec:app}
The application is built with Streamlit and Plotly and reads only from the database and from the result files
of the other components. It has six pages (Table~\ref{tab:pages}).

\begin{table}[ht]
\centering
\small
\begin{tabular}{@{}lp{12cm}@{}}
\toprule
Page & Content \\
\midrule
Home & Corpus size, number of events, documents per year and source type, largest sources \\
Corpus & Full-text search with filters on source type, source and year range; full text and link for each result \\
Timeline & Events by category over time, events per month, filters (category, organisation, capability, autonomy
level, origin, months), event details with source links \\
Phases & Events per month with a three-month moving average, selectable feature trends, dominant category per
month, the quarterly 2017--2021 view and change-point segments \\
Forecast & Predicted dominant category for the next month, class probabilities, baselines, test metrics,
validation and test history, top features and supporting events \\
Evidence & One event with its fields, its source documents and the trigger phrase shown in its surrounding text \\
\bottomrule
\end{tabular}
\caption{Pages of the application.}
\label{tab:pages}
\end{table}

Four design decisions follow from the project's emphasis on honest interpretation.
\begin{enumerate}
\item \emph{Gold is the default.} The timeline opens on the hand-labelled events. Extracted events can be added
with the origin filter, and whenever they are included the page shows a warning that they are model output with
rough categories and approximate counts.
\item \emph{Traceability.} From any forecast, supporting event or timeline point the user can open the Evidence
page and read the original document and the exact words that signalled the event.
\item \emph{No invented results.} The forecast page reads the forecasting output from a result file. If that
file is missing, the page shows a clearly labelled most-frequent baseline instead of a model forecast. The
history is split into validation and test months so that the displayed accuracy matches the reported test
metric.
\item \emph{Empty states.} Pages explain what is missing (for example a missing database or result file) and
which command creates it, instead of failing.
\end{enumerate}

%------------------------------------------------------------------------------
\subsection{Reproducibility and engineering practice}
\label{sec:repro}
\begin{itemize}
\item \emph{Commands.} Collection, the cleaning pipeline, the database build, the loading of extracted events and
the application each have one command, listed in the repository README together with an installation guide and
a demo script.
\item \emph{Configuration.} Sources, the taxonomy (categories, phases, scales and tag lists) and the
forecasting split and cutoff are YAML files rather than constants in code, so all components share one definition.
\item \emph{Fixed inputs.} The corpus is identified by a SHA-256 value; package versions are pinned in the
requirements file.
\item \emph{Tests.} Unit tests cover the storage module (schema, loaders, search, validation, event queries and
the gold merges) and the cleaning and date logic; the full suite of 32 tests passes.
\item \emph{Version control.} Each member worked on a branch named after their initials and task, changes were
merged through reviewed pull requests, the main branch was protected, and raw data, databases, keys and large
files were kept out of the repository.
\end{itemize}

%------------------------------------------------------------------------------
\subsection{Limitations of the data infrastructure}
\label{sec:limits}
\begin{itemize}
\item \emph{Coverage.} The years 2017--2021 and 2022 contain fewer documents than planned (58--85 per year before
2022 and 174 in 2022), so the early part of the timeline is sparse. No Meta, Mistral, DeepSeek or Qwen sources
were collected, and Microsoft is represented by only ten documents. The news and incident-report families
planned for the project are not part of the frozen corpus.
\item \emph{Source bias.} The corpus favours a handful of laboratories and frameworks. The monthly caps reduce,
but do not remove, the effect of prolific sources, and caps themselves influence monthly counts.
\item \emph{Dates.} For papers and release notes the publication date is used as the event date because the
text rarely states a separate date, so the timeline can lag the real development.
\item \emph{Text quality.} Some blog posts are available only as feed summaries, which gives the extraction less
text to work with.
\item \emph{Extracted events are provisional.} The 871 extracted events use probability thresholds chosen on
about 37 labelled documents per source type. The set contains no release-note events, its categories are about
60\% accurate against the gold events, and its counts are approximate. It is displayed but not used for
forecasting.
\item \emph{Deployment.} The application is a local, single-user analysis tool over a read-only database, not a
production service.
\end{itemize}

\end{document}
